%Knoten- und Maschenanalyse
\newpage

\newvideofile{KnotMasch}{Knot- and Mesh analysis}
\begin{frame}
	\fta{Knot- and Mesh analysis}

	\s{When considering electrical networks, the currents 
    in nodes and the voltages in meshes are analysed in particular. The analysis of nodes and 
    meshes in electrical networks is carried out using the two Kirchhoff's rules.}

    \b{Die Analyse von Knoten und Maschen in elektrischen Netzwerken erfolgt durch 
    die Kirchhoffschen Regeln.}

    \begin{Lernziele}{Knot- and Mesh analysis}
        The Students
        \begin{enumerate} 
            \item know the basic definitions for describing an electrical network.
            \item can apply the node rule and the mesh rule to electrical networks. 
        \end{enumerate}
    \end{Lernziele}
    \speech{Knotmasch}{1}{
        Nodes and meshes are definitions that enable the analysis of electrical networks. 
        The nodes and meshes of an electrical network are treated using 
        Kirchhoff's rules. Here we will explain the basic 
        definitions of an electrical network and how the nodes and meshes are analysed using 
        Kirchhoff's rules.}
\end{frame}


%Herleitung Knotenregel
\begin{frame}
    \renewcommand{\figurename}{Figure}
    \ftb{Current density in free space} \label{Stromdichte}
 
    \s{
        Facts and concepts from field theory lead to explanations in
        structured field space. To this end, Figure \ref{BildStromdichte}
        shows a positively charged plate and a negatively charged plate
        facing each other.The space between the two plates has an
        electric field $\vec{E}$ and a medium with conductivity $\sigma$. Charge carriers
        can move in the free space. A directed movement
        of charge carriers and thus a current density $\vec{J}$ between the plates occurs.
        The current density is source-free in the area under consideration.
    }

    \b{Electrical Field $\vec{E}$ between two charged plates:
    \begin{itemize}
        \item<1->   Raum zwischen den Platten ist mit einem leitfähigen Material $\sigma$ gefüllt
        \item<2->   Stromdichte $\vec{J}$ im elektrischen Feld 
        \item<3->   Quellenfreiheit der Stromdichte (Kontinuitätsgleichung)
    \end{itemize}}

    \fu{
        \input{Tikz/src/elektrischesfeld.tex}
        }{{\bf Electric field between a positively charged and a negatively charged 
        plate.} The current density of a sample area in the electric field is source-free.
        \label{BildStromdichte}}

    \s{
        The source-free nature of the current density is illustrated by the equation \ref{divJQuellenfreiheit}.
        This describes that no sources or sinks exist, as there are no field-forming charges in space.
        In this case, the current density is source-free and the same number of charges flow into and out of the surface.
        Mathematically, this is expressed in the continuity equation for consideration in free space:
    }

    \only<3->{
    \begin{eq}
        \oint_{}^{} \vec{J} d\vec{A} = 0 \quad \rightarrow \quad div \vec{J} = 0 \label{divJQuellenfreiheit}
    \end{eq}
    }
    
    \speech{StromdichteimfreienRaum}{1}{
        Zuvor jedoch noch ein paar Grundlagen. Es wird ein elektrisches Feld E, zwischen einer positiv geladenen Platte, 
        hier in rot, und einer negativ geladenen Platte, hier in blau, betrachtet. Der Raum zwischen den beiden geladenen
        Platten ist mit einem leitfähigen Material Sigma, gefüllt und geometrisch nicht näher definiert.}
    \speech{StromdichteimfreienRaum}{2}{
        Aus dem elektrischen Feld E, resultiert eine theoretische Stromdichte J, im Raum zwischen den geladenen Platten.}
    \speech{StromdichteimfreienRaum}{3}{
        Die Kontinuitätsgleichung beschreibt das Oberflächenintegral der Stromdichte J, entlang einer geschlossenen Linie, 
        welche die Fläche A, einschließt, also dass die Stromdichte an jedem Punkt gleich ist und nicht zu- oder abnimmt. 
        Hiermit wird erklärt, dass die Divergenz der Stromdichte bei stetigen Strömen Quellenfrei beziehungsweise frei von 
        Senken ist.}
\end{frame}


\begin{frame}
    \renewcommand{\figurename}{Figure}
	\ftb{Current density in structured space}

    \s{
        In a structured field space (see Figure \ref{BildLeiterknoten}), the current density must also be free of sources.
        All charges that enter the field space must also leave it again.
        This also applies in cases where an electrical conductor has more than one discharge path
        . A supply line with two conductors 
        is indicated, in which the charges are distributed between the two conductors. The cylinder surface areas
        $A_1$, $A_2$ and $A_3$ define a structured field space with the adjacent cylinders. Just as the charges are distributed, 
        the current densities are distributed throughout the structured space. Thus, the current density
        $\vec{J_1}$ is the sum of the outgoing current densities $\vec{J_2}$ and $\vec{J_3}$. 
        
        \begin{columns}
        \column[c]{0.9\textwidth}  
        \fu{
            \input{Tikz/src/strukturierterfeldraum.tex}
        }{{\bf Structured field space.} Simplified representation of a splitting conductor with current densities 
        $J_1$, $J_2$ and $J_3$ and areas $A_1$, $A_2$ and $A_3$.\label{BildLeiterknoten}}    
        \end{columns}

        Since the field space in the example is structured by the surfaces and the cross-section of the 
        conductor is defined, the currents of the conductors can be determined from the relationship between the current density and the 
        surface area according to equation \ref{GleichungStromberechnung}. 

        \begin{eq}
            I_1 = \iint_{A}^{} \vec{J_1} d\vec{A}  \label{GleichungStromberechnung}
        \end{eq}

        The relationship between the current densities 
        can thus be transferred to the currents. The current of the supply line $I_1$ is the sum of the
        two branch currents $I_2$ and $I_3$. The example conductor can also be represented as a simplified 
        electrical network. The currents determined around the connection point K (node) 
        are shown in Figure \ref{FigureNodeNetwork}. 

        \fu{
            \input{Tikz/src/knotenstrukturiert.tex}
        }{{\bf Example nodes from the structured space.} The current densities $J_1$, $J_2$ and $J_3$ result in the currents 
        $I_1$, $I_2$ and $I_3$ around node K. \label{BildKnotennetzwerk}}
    }

    \b{
    
    Definierter Feldraum:

    \begin{minipage}[t]{0.7\textwidth}
        \begin{itemize}
            \item<2-> Bekannte Flächen und Stromdichten
            \item<3-> Strom
            \item<4-> Netzwerk
        \end{itemize}
    \end{minipage}
    \begin{minipage}[t]{0.29\textwidth}
        \only<3->{
        \begin{eq}
            I_1 = \iint_{A}^{} \vec{J_1} d\vec{A} 
        \end{eq}
        }
    \end{minipage}
    

    
    \begin{minipage}{1\textwidth}			
        \resizebox{\textwidth}{!}{\centering{
            \begin{circuitikz}
                \draw
                [line width = 2.5, black]
                (0,0) circle (28pt)
                (0,1) -- (10,3)
                (0,-1) -- (5,0)
                (7,0.4) -- (10,1)
                (5,0) -- (5,-2.6)
                (7,0.4) -- (7,-2.6)
                (10,1) arc (-90:90:28.5pt)
                (5,-2.6) arc (-180:0:28.5pt);
                \pause
                
                \draw[line width = 1.5, red][double,->] (1.5,0.3) -- (2.5,0.5);
                \draw[red](1.5,0.3) node[left] {$J_1$};
                \draw[line width = 1.5, red][double,->] (9,1.8) -- (10,2);
                \draw[red](10,2) node[right] {$J_2$};
                \draw[line width = 1.5, red][double,->] (6,-1.6) -- (6,-2.6);
                \draw[red](6,-2.6) node[below] {$J_3$};
                \draw
                [line width = 1.5, blue]
                (3,0.6) ellipse (14pt and 28pt)
                node at (3,0.6) {$A_1$}
                (8.5,1.7) ellipse (14pt and 28pt)
                node at (8.5,1.7) {$A_2$}
                (6,-1.1) ellipse (28pt and 14pt)
                node at (6,-1.1) {$A_3$};
                \pause

                \draw[line width = 1.5, red][double,->] (5,1) -- (7,1.4);
                \draw[red](5,1) node[left] {$I_1$};
                \draw[red](7,1.4) node[right] {$I_2$};
                \draw[line width = 1.5, red][double,->] (6,1.2) -- (6,0.2);
                \draw[red](6,0.2) node[below] {$I_3$};
                \pause

                \draw
                (15,0) to [short, i=\red{$I_1$},-*] (17,0)
                (17,0) to [short, i=\red{$I_2$}] (19,0)
                (17,0) to [short, i=\red{$I_3$}] (17,-2);
            \end{circuitikz}	
        }}	
    \end{minipage}
	} 
    
    \speech{StromdichteimstrukturiertenRaum}{1}{
        Zuvor haben wir die Stromdichte im freien Raum betrachtet. Jetzt soll die Stromdichte in einem strukturierten Raum,
        zum Beispiel in einem elektrischen Leiter, genauer Betrachtet werden.}
    \speech{StromdichteimstrukturiertenRaum}{2}{
        Im strukturierten Raum sind die Flächen des Raumes und die Stromdichten bekannt. Hier werden die Stromdichten J-Eins,
        J-Zwei und J-Drei mit ihren zugehörigen Flächen markiert.}
    \speech{StromdichteimstrukturiertenRaum}{3}{
        Aus den Stromdichten und den zugehörigen Flächen lassen sich die daraus resultierenden elektrischen Ströme bestimmen.
        Hier teilt sich der Strom I-Eins auf die Ströme I-Zwei und I-Drei auf.}
    \speech{StromdichteimstrukturiertenRaum}{4}{
        Die dreidimensionale Abzweigung im Leiter lässt sich auch als zweidimensionales elektrisches Netzwerk abbilden.}
\end{frame}

%Exkurs: Graphentheorie
\begin{frame}
    \ftb{Digression: Graph theory}
    
    \s{
        In mathematics, graph theory deals with the description of
        nodes and edges, while account analysis also defines nodes and branches. 
        Graph theory allows the edges and nodes shown in Figure 
        \ref{GraphentheorieKnoten} and their properties 
        to be related to each other. Multiple nodes can be connected to each other via the edges. 
    }

    \b{
        Beziehung zwischen Knoten und Ästen in der Graphentheorie:
        \begin{itemize}
            \item<2-> Kanten $\rightarrow$ Äste
            \item<3-> Knoten
        \end{itemize}
    }

    \begin{columns}
        \renewcommand{\figurename}{Figure}
        \column[c]{0.3\textwidth}    
            \fu{
                \input{Tikz/src/graphentheorieknoten.tex}
            }{{\bf Example nodes in graph theory.} Node with 4 edges to explain nodes and edges in graph theory.
            \label{GraphentheorieKnoten}}      
    \end{columns}
        
    \speech{KantenKnoten}{1}{
        Um elektrische Netzwerke einfacher beschreiben zu können, werden Bestandteile der Graphentheorie hinzugezogen.
        Die Graphentheorie wird in der Elektrotechnik zur Analyse von elektrischen Netzwerken und der Analyse von 
        Versorgungsnetzen eingesetzt.} 
    \speech{KantenKnoten}{2}{    
        Dafür werden Kanten, oder auch Äste definiert. Sie dienen als Verbindungswege zwischen zwei Punkten.}
    \speech{KantenKnoten}{3}{    
        Die Kanten können an Knoten zusammengeführt werden. Hier in rot markiert.}
\end{frame}

%Netzwerk und Matrix Graphentheorie
\begin{frame}
    \renewcommand{\figurename}{Figure}
    \ftx{Exkurs: Graphentheorie}

    \s{
        When several nodes are connected to each other with edges, a graph is created. This 
        graph can be used to describe an electrical network, for example. 
        The example of a network with five nodes and the edges between the nodes
        is shown in Figure \ref{GraphentheorieNetzwerk}. When searching for a path through the network
        , there are various possibilities. A path from node 1 via nodes 2,
        4 and 5 back to node 1 is highlighted in red. If the start and end are the same node, 
        this path is also referred to as a cycle in graph theory. Any cycle in which each node 
        is "visited" only once, except for the start and end, is called a "circle". In 
        electrical networks, these circles are also referred to as meshes. 
    }

    \b{
        Graph:
        \begin{itemize}
            \item<2-> Weg
            \item<3-> Weganfang = Wegende $\rightarrow$ Zyklus
            \item<4-> Kein kreuzender Zyklus $\rightarrow$ Kreis
        \end{itemize}
    }

   
    \fu{
        \input{Tikz/src/graphenteorienetzwerk.tex}            
    }{{\bf Example network from graph theory.}
    Network with 5 nodes and edges running between the nodes to 
    explain a cycle in graph theory. \label{GraphentheorieNetzwerk}}      


    \s{
        In mathematics, a matrix is a useful tool for explaining the relationship between nodes. 
        The nodes have a differentiated number of edges, and not all nodes are 
        connected to each other. The adjacency matrix describes the relationship between the nodes. 
        Such an adjacency matrix is set up in equation \ref{AdjazenzMatrix} for the
        example network described. It thus completely describes the graph, apart from the graphical 
        representation in the 2D display. 
    }

    \b{
        \onslide<5->{
            Zugehörige Adjazenz-Matrix:
        }
    }

    \onslide<5->{
    \begin{eq}
        A = 
        \begin{bmatrix}
            \color{blue}{1} & 1 & 0 & 0 & 1\\
            1 & \color{blue}{1} & 1 & 1 & 1\\
            0 & 1 & \color{blue}{1} & 1 & 0\\
            0 & 1 & 1 & \color{blue}{1} & 1\\
            1 & 1 & 0 & 1 & \color{blue}{1}
        \end{bmatrix} \label{AdjazenzMatrix}
    \end{eq}
    }

    \speech{Zyklus}{1}{
        Werden mehrere Kanten und Knoten verschaltet, ergibt sich daraus ein sogennanter Graf.} 
    \speech{Zyklus}{2}{    
        In diesem Graf lässt sich von einem beliebigen Knoten zu einem anderen Knoten ein Weg definieren.}
    \speech{Zyklus}{3}{    
        Ist der Knoten über andere Knoten hinweg sowohl Weg-Anfang als auch Weg-Ende ergibt sich daraus ein Zyklus.}
    \speech{Zyklus}{4}{    
        Wenn sich der Weg in einem Zyklus nicht kreuzt ergibt sich ein Kreis. Diese Kreise werden später bei der Analyse von 
        Maschen noch eine Rolle spielen.}
    \speech{Zyklus}{5}{    
        Die Beziehungen zwischen Knoten und Kanten in einem Graf können auch mathematisch durch die AdJazenz-Matrix beschrieben 
        werden. In die AdJazenz-Matrix wird eine Eins eingetragen, wenn ein Knoten über eine Kante einen direkten Weg zu einem 
        anderen Knoten hat. Für den Fall, dass keine direkte Verbindung zwischen zwei Knoten über eine Kante besteht, wird eine 
        Null notiert.}
\end{frame}

\begin{frame}
    \ftx{Graphentheorie}

    \begin{Merksatz}{Graphentheorie}
        In graph theory, networks are described using edges and nodes. When analysing networks, paths and cycles are defined.
        A continuous cycle with an identical starting point and end point is referred to as a circle. 
    \end{Merksatz}

    \speech{Graphentheorie}{1}{
        In graph theory, networks are described using edges and nodes. When analysing networks, paths and cycles are defined.
        A continuous cycle with an identical starting point and end point is referred to as a circle.} 
\end{frame}


\begin{frame} 
    \renewcommand{\figurename}{Figure}
	\ftb{Knot, Branch, Mesh}

	\s{
		An electrical network consists of nodes, branches and meshes. A branch connects exactly two nodes by means of 
        one or more circuit elements. The same current flows through all elements of a branch. 
        There are three branches in Figure \ref{Netz1}. One branch with the components $R_1$, $R_5$ and $U_\mathrm{q1}$, one 
        branch with the resistance $R_4$ and one branch with the components $R_2$, $R_3$ and $U_\mathrm{q2}$.
        A node is a point at which at least three connections of the circuit elements converge. The current 
        can therefore split here. The electrical potential is identical for all connected terminals.
 
        In the circuit diagram, a node is indicated by a filled circle. However, if two or more 
        of these circles are connected to each other by only one line, it is a single node, since the 
        potential is also the same here. A node is denoted by $K_n$.

		A mesh is a closed path consisting of at least two branches. In the net in Figure 
        \ref{Netz1}, three meshes can be defined:
			\begin{itemize}
				\item $M_1$ consisting of $R_1$, $R_4$, $R_5$ and $U_\mathrm{q1}$
				\item $M_2$ consisting of $R_2$, $R_3$, $R_4$ and $U_\mathrm{q2}$
				\item $M_3$ consisting of $R_1$, $R_2$, $R_3$, $U_\mathrm{q2}$, $R_5$ und $U_\mathrm{q1}$
			\end{itemize}
		A mesh is designated by $M_\mathrm{n}$. The direction of rotation is important and is indicated by an arrow.
	}

	\fu{	
        \input{Tikz/src/Netz1.tex}
	}{{\bf Direct current network. }Network with {\color{red}Knots}, Branches and {\color{blue}Meshes}\label{Netz1}}

    \speech{Knoten,Zweige,Maschen}{1}{ Im folgenden werden die Definitionen von Knoten, Zweigen und Maschen in elektrischen Netzwerken 
        näher erläutert. Hierfür wird das dargestellte elektrische Netzwerk genutzt. Wie schon im Exkurs zur Graphentheorie ausgeführt, 
        sind Kanten oder Äste Verbindungswege zwischen zwei Punkten. Diese Kanten werden in der Elektrotechnik als Zweige bezeichnet. 
        Dieses elektrische Netzwerk setzt sich aus insgesamt drei Zweigen zusammen .} 
    \speech{Knoten,Zweige,Maschen}{2}{ Die Punkte, welche in dem gezeigten Netzwerk verbunden werden, sind die Knoten Ka 1 und Ka 2 .} 
    \speech{Knoten,Zweige,Maschen}{3}{ Eine Masche beschreibt einen geschlossenen Weg, der aus mindestens zwei Zweigen besteht. In dem 
        Netzwerk können so beispielsweise die Masche M 1 über die Komponenten R 1, R 4, R 5 und U Q 1 sowie die Masche M 2 über die 
        Komponenten R 2, R 3, U Q 2 und R 4 bestimmt werden. Außerdem können Maschen auch über zweige hinaus andere Zweige kreuzen. So 
        kann in diesem elektrischen Netzwerk die Masche M 3 über die Komponenten R 1, R 2, R 3, U Q 2, R 5 und U Q 1 definiert werden .} 
    \speech{Knoten,Zweige,Maschen}{4}{Bei den Maschen ist die Umlaufrichtung von Bedeutung und wird mit einem Pfeil gekennzeichnet. ebenso
        werden die richtungen von elektrischen Strömen gekennzeichnet. Hierzu später mehr bei der Knotenregel und der Maschenregel. } 
\end{frame}	

\begin{frame}
    \ftx{Knots, Branches and Meshes}

    \begin{Merksatz}{Knots, Branches and Meshes}
        Electrical networks are described by nodes, branches and meshes, analogous to graph theory.
    \end{Merksatz}

    \speech{MerkeKZM}{1}{Elektrische Netzwerke werden analog zur Graphentheorie durch Knoten, Zweige und Maschen beschrieben. } 
\end{frame}



\begin{frame} 
    \renewcommand{\figurename}{Figure}
	\ftb{The complete tree}
	\s{
		For further analysis of the network, it is necessary to determine the network equations. However, this linear 
        system of equations is always overdetermined, which is why it is necessary to reduce the number of equations. 
        It is important to identify and eliminate the linearly dependent equations.
                
        The nodes, branches and meshes of a network can be represented in a graph that only shows the 
        connections between them. To represent the graph, all branches are shown as lines that 
        connect the nodes. The content of the branch is irrelevant for this purpose. Figure \ref{AbbGraph} shows an 
        example network with the resulting graph.}

	\fu{
		\begin{minipage}{0.4\textwidth}
			\b{\resizebox{\textwidth}{!}{\centering{
				\begin{circuitikz}
					\draw (0,0) to[short] (0,1.5) 
						to[V,v<, name=U0, *-*] (0,4.5)
						to[short] (0,6)
						to[V, name=U1] (3,6)
						to[short] (3,4.5)
						to[R=$R_3$, *-*](3,1.5)
						to[R=$R_4$](0,1.5);
					\draw (0,4.5) to[R=$R_2$] (3,1.5);
					\draw (0,4.5) to[R=$R_1$] (3,4.5)
						to[short] (4,4.5)
						to[short] (4,0)
						to[short] (3,0)
						to[I, name=I2] (0,0);
					\varrmore{U0}{$U_{0}$};
					\varrmore{U1}{$U_{1}$};
					\iarrmore{I2}{$I_{2}$};
				\end{circuitikz}
			}}}		
			\s{
				\input{Tikz/src/Gleichstromnetzwerk.tex}
			}
		\end{minipage}\hfill\pause
		\begin{minipage}{0.26\textwidth}
			\resizebox{\textwidth}{!}{
				\input{Tikz/src/Graph.tex}
			}
		\end{minipage}	
	}{{\bf Direct current network and graph.} Generation of the graph from a network to be calculated\label{AbbGraph}}

    \speech{VollstandigeBaum}{1}{Bei der Analyse von elektrischen Netzwerken kann der vollständige Baum genutzt werden.}
    \speech{VollstandigeBaum}{2}{Der vollständige Baum erklärt dabei ein elektrisches Netzwerk als ein gebilde aus Zweigen und Knoten.}
         
\end{frame}

\begin{frame} 
    \renewcommand{\figurename}{Figure}
	\ftx{The complete tree}\label{vollstbaum}
	
	\s{
		In this graph, a complete tree can be drawn that contains all nodes but does not itself form a mesh. 
        The tree branches connect all nodes to each other but do not form a closed line. 
        The tree branches can form branches, meaning that a node can also touch more than two tree branches. In 
        larger networks, there are several possibilities for a complete tree, which are in principle equivalent.
    However, when presenting the various calculation methods, certain combinations are preferred depending on the method.
    A complete tree contains exactly $k-1$ branches.
		}
	
	\fu{
		\begin{minipage}{0.19\textwidth}
			\input{Tikz/src/bspbaum1.tex}
		\end{minipage}
		\hspace{0.2cm}
		\begin{minipage}{0.19\textwidth}
			\input{Tikz/src/bspbaum2.tex}
		\end{minipage}
		\hspace{0.2cm}
		\begin{minipage}{0.19\textwidth}
			\input{Tikz/src/bspbaum3.tex}
		\end{minipage}	
		\hspace{0.2cm}
		\begin{minipage}{0.19\textwidth}
			\input{Tikz/src/bspbaum4.tex}
		\end{minipage}
	}{{\bf Example trees.} Examples of different complete trees of a network drawn in red}
	
	\b{
		\begin{itemize}
			\item<2-> \red{Baumzweige} $\rightarrow$ k-1 des vollständigen Baumes
			\item<2-> die übrigen Zweige sind Verbindungszweige
		\end{itemize}
	}

	\s{
		The branches belonging to the tree are called tree branches (red), the other connecting branches (black).
        To find the linearly independent meshes, a mesh is formed for each connecting branch that contains only tree branches except for 
        the connecting branch.
		}
    \begin{itemize}
        \item<3-> A network with $z$ branches and $k$ nodes contains $z-k+1$ connecting branches.
    \end{itemize}
	
    \speech{VollstandigeBaumBsp}{1}{Beim vollständigen Baum gibt es eine vielzahl an möglichen Anordnungen von verwendeten Zweigen.}
    \speech{VollstandigeBaumBsp}{2}{Hierbei muss zwischen den hier in rot gestrichelten Baumzweigen und den Verbindungszweigen 
    unterschieden werden. Es gibt in einem elektrischen Netzwerk immer k minus eins Baumzweige. Über die Baumzweige werden alle Knoten
    miteinander verbunden, ohne dabei eine Masche zu definieren, also ohne, dass ein Knoten zwei mal berührt wird. 
    Die übrigen Zweige werden als Verbindungszweige definiert. }
    \speech{VollstandigeBaumBsp}{3}{Ein Netzwerk mit Z Zweigen und K Knoten enthält Z minus K plus eins Verbindungszweige.
    Die Anzahl der Verbindungszweige entspricht später auch der Anzahl der unabhängigen Maschengleichungen, welche nötig sind, 
    um ein elektrisches Netzwerk zu analysieren.}

\end{frame}

\begin{frame}
    \ftx{The complete tree}

    \begin{Merksatz}{The complete tree}
        With the help of a complete tree, all nodes are connected to each other by tree branches.
        This always results in k - 1 tree branches.
        The complete tree does not form a mesh. 
    \end{Merksatz}

    \speech{VollstandigeBaumMrk}{1}{Mithilfe eines vollständigen Baumes werden alle Knoten miteinander durch Baumzweige verbunden.
    Es ergeben sich dabei immer k minus 1 Baumzweige. Der vollständige Baum bildet keine Masche.}
\end{frame}


%Knotenregel
%Knotenregel Summenzeichnen
\begin{frame}
    \renewcommand{\figurename}{Figure}
    \ftb{Knot rule}

    \s{
        If more than two lines meet at a point in an electrical 
        network, this is referred to as a node K. The direction of the arrow from 
        the perspective of the node determines the sign of the current. If a current 
        flows into a node, the current arrow points towards the node and this is 
        referred to as an inflow. If the current arrow points out of the node, 
        the current flows out of the node and this is an outflow. 
        The node analysis is based on Kirchhoff's first law, which states:
        The sum of all currents at a node is zero. (Gleichung \ref{Knotenregel}). 
        This law follows directly from the continuity equation. (vgl. Abschnitt 
        \nameref{Stromdichte})
    }

    \b{ 
        \begin{itemize}
            \item Die Summe aller Ströme an einem Knoten ist gleich Null.
        \end{itemize}
    }

    \begin{eq}
        \sum_{k = 1}^{N}{I_\mathrm{k} = 0} \label{Knotenregel}
    \end{eq}
    
    \s{
        Thus, for each node, the sum of the currents from inflows 
        and outflows must balance each other out. Figure \ref{BildBeispielknoten} shows a 
        node K and the inflows and outflows. The flows $I_1$ and 
        $I_2$ are inflows. The three flows $I_3$, $I_4$ and $I_5$ are 
        outflows.
 
}

    \b{
        \begin{itemize}
            \item Beispielnetzwerk:
        \end{itemize}
    }

        \fu{
            \input{Tikz/src/beispielknoten.tex}
        }{{\bf Example node.} Node for node analysis with two incoming and 
        three outgoing branches. \label{BildBeispielknoten}}


	\s{
		If the currents are applied one after the other and set to zero, 
        the equation \ref{KnotenregelBeispielknoten} is obtained. Depending on whether 
        it is an inflow or an outflow, the sign is selected. 
        Inflows are assigned a positive sign and outflows a 
        negative sign.
 
	}
    
    \b{
        \begin{itemize}
            \item Summe aus Einströmungen und Ausströmungen gleich null:
        \end{itemize}
    }
    
    \begin{eq}
        \sum_{k = 1}^{N}{I_\mathrm{k}} =0=I_1+I_2-I_3-I_4-I_5    \label{KnotenregelBeispielknoten}
    \end{eq}

    \speech{Knotenregel}{1}{Die Knotenregel beschreibt die Verhältnisse von elektrischen Strömen an einem Knoten und besagt, 
    dass die Summe aller Ströme an einem Knoten gleich null ist.
    Anhand des Beispielnetzwerkes mit den fünf Strömen soll die Knotenregel veranschaulicht werden. Bei den Strömen I 1 und I 2 handelt 
    es sich um elektrische Ströme die auf den Knoten zuströmen. Die Ströme I 3, I 4 und I 5 strömen von dem Knoten weg. 
    Sie Summe aus Einströmungen und Ausströmungen muss also gleich null sein.}
\end{frame}

\begin{frame}
    \ftx{Knotenregel}

    \begin{Merksatz}{Knot rule}
        The sum of the incoming flows is equal to the sum of the outgoing flows, or the sum of all flows at a node
        is zero:
        \begin{eq}
            \sum_{k = 1}^{N}{I_\mathrm{k} = 0} 
        \end{eq}
    \end{Merksatz}

    \speech{KnotenregelMrk}{1}{Die Summe der zufließenden Ströme ist gleich der Summe der abfließenden Ströme oder die Summe aller Ströme 
    an einem Knoten ist Null.}
\end{frame}



%Maschenanalyse
\begin{frame}
    \ftb{Mesh rule}

    \s{
		When a particle moves between two points in an electric field, work is done.
        If the starting point is identical to the end point, the same amount of work must be delivered as absorbed.
        In this case, the sum of the work is zero. If such a complete mesh cycle is performed along a
        closed path, the sum of the partial voltages must also be zero.         
        This assumes that there is no time-varying magnetic flux density and therefore no 
        electrical eddy field is created (see equation \ref{GleichungRingintegralFeldstärke}). The equation
        describes Faraday's law of induction for static fields (here: no change in a magnetic field).
    }

    \b{
        Ringintegral über der Feldstärke: 
        \begin{itemize}
            \item Bedingung: Keine zeitlich veränderte magnetische flussdichte und somit kein elektrisches Wirbelfeld. 
        \end{itemize}
    }

    \begin{eq}
        \oint_{s}^{} \vec{E} d\vec{s} = 0 \quad \rightarrow \quad rot \vec{E} = 0 \label{GleichungRingintegralFeldstärke}
    \end{eq}

    \s{
        According to Kirchhoff's second law, mesh analysis states that for a 
        complete circuit (mesh) in an electrical network, the sum of all 
        voltages is equal to zero. This is illustrated by the equation \ref{Maschenregel}. 
    }

    \b{\onslide<2->{
        Maschenregel:
        \begin{itemize}
            \item Bei einem vollständigen Umlauf in einem elektrischen Netzwerk ist die Summe aller Spannungen gleich null. 
        \end{itemize}}
    }

    \onslide<2->{
    \begin{eq}
        \sum_{k = 1}^{N}{U_\mathrm{k}} =0    \label{Maschenregel}
    \end{eq}}

    \speech{Maschenregel}{1}{Tritt keine zeitliche Änderung bei elektrischen und magnetischen Feldgrößen auf, können keine 
    elektrischen Wirbelströme induziert werden. Bewegt sich unter dieser Bedingung ein Teilchen zwischen zwei Punkten im 
    elektrischen Feld, wird Arbeit verrichtet. Ist der Startpunkt identisch mit dem Endpunkt, muss ebenso viel Arbeit 
    abgegeben wie aufgenommen werden. Die Summe der Arbeit ist in diesem Fall gleich Null. Ähnlich zum Umlauf im freien Raum,
    wie dem elektrischen Feld, verhält sich ein Umlauf im strukturierten Raum, also in einem elektrischen Netzwerk.}
    \speech{Maschenregel}{2}{Hierzu wird die Maschenregel definiert. Die Maschenregel besagt, dass die Summe aller Teilspannungen 
    bei einem vollständigen Umlauf gleich null sein muss.}
\end{frame}

\begin{frame}
    \renewcommand{\figurename}{Figure}
    \ftx{Maschenregel}
    
    \s{
        Within a mesh, voltages with an arrow direction in the same direction are counted as positive,
        while voltages with an arrow direction opposite to the direction of rotation are counted as negative.
        Based on the current curves and potentials from Figure \ref{BeispielnetzwerkMaschenanalyse}, 
        the voltages of the respective branches can be determined. 
        There is no potential difference above the parallel connection of the resistors. 
        The same applies to the electrical connection below the parallel connection. The current $I_\mathrm{R1}$ 
        flows through the resistor $R_\mathrm{1}$. According to Ohm's law, this results in the voltage $U_\mathrm{R1}$ 
        across $R_\mathrm{1}$. Equivalently, the voltage $U_\mathrm{R2}$ across $R_\mathrm{2}$ is determined.
    }

	\fu{
		\input{Tikz/src/beispielnetzwerkmaschenanalyse.tex}
	}{{\bf Example network with one power source and two parallel resistors.} In mesh $M_1$, the voltages
    $U_\mathrm{R1}$ and $U_\mathrm{R2}$ must be identical and cancel each other out according to their directions in the mesh. \label{BeispielnetzwerkMaschenanalyse}}

    \s{
        According to the network presented and the direction of the drawn mesh $M_\mathrm{1}$,
        the relationship according to equation \ref{Maschengleichung} applies. The voltage $U_\mathrm{R2}$ runs 
        in the same direction as the drawn mesh and is assigned a positive sign . 
        The voltage $U_\mathrm{R1}$ runs in the opposite direction to the mesh $M_\mathrm{1}$ and is therefore negative. 
        In total, these two voltages must add up to zero. 
    }

    \onslide<2->{
    \begin{eq}
        \sum_\mathrm{k = 1}^{N}{U_\mathrm{k}} \onslide<3->{= + U_\mathrm{R2}} \onslide<4->{- U_\mathrm{R1}} \onslide<5->{=0}  \label{Maschengleichung}
    \end{eq}}

    \speech{MaschenregelErkl}{1}{Zur Veranschaulichung der Maschenregel wird die Masche M 1 im gezeigten Netzwerk erklärt.}
    \speech{MaschenregelErkl}{2}{Der Maschenumlauf der Masche M 1 verläuft im Uhrzeigersinn. Nach der Maschenregel muss die 
    Summe der Teilspannungen Null ergeben.}
    \speech{MaschenregelErkl}{3}{Hierbei wird die Richtung der Masche und der Teilspannungen ausschlaggebend. Die Richtung 
    der Teilspannung U R 2 über den Widerstand R 2 verläuft in derselben Richtung wie die Umlaufrichtung der Masche M Eins. Diese 
    Teilspannung wird deshalb mit einem positiven Vorzeichen versehen.}
    \speech{MaschenregelErkl}{4}{Die Richtung der Teilspannung U R 1 über dem Widerstand R 1 ist der Umlaufrichtung der 
    Masche M 1 entgegengesetzt und erhält somit ein negatives vorzeichen.}
    \speech{MaschenregelErkl}{5}{Unter der Annahme, dass die beide parallel zueinander geschalteten Widerstände R 1 und 
    R 2 betragsmäßig über eine identische Teilspannung verfügen, ergibt sich so bei einem Maschenumlauf eine Gesamtspannung 
    von Null.}
\end{frame}


\begin{frame}
    \ftx{Maschenregel}

    \begin{Merksatz}{Mesh rule}
        In a mesh, the sum of all partial stresses is equal to zero.:
        \begin{eq}
            \sum_{k = 1}^{N}{U_k} =0
        \end{eq}
    \end{Merksatz}

    \speech{MaschenregelMrk}{1}{In einer Masche ist die Summe aller Teilspannungen gleich Null}
\end{frame}



\begin{frame} \ftb{Knot and Mesh analysis}
	\s{In node and mesh analysis, the complete system of equations consisting of node and mesh equations is set up and then solved. The procedure is simple, but generates a number of equations for a network that corresponds to the number of branches. The procedure is as follows:}
	\b{Ziel: Berechnung eines vollständigen Netzwerkes.
	Arbeitsschritte:}
	
	\begin{enumerate}
		\item<2-> Simplifying the network.
		\item<3-> Draw all current arrows and create the graph.
		\item<4-> Set up $k-1$ node equations. \s{Any node equation can be omitted, since in a network with $k$ nodes only $k-1$ node equations are linearly independent of each other. In principle, it does not matter which node equation is omitted. It is best to omit the most complicated equation.}
		\item<5-> Set up the $m=z-k+1$ linearly independent mesh equations. \s{The equations are created using the complete tree.}
		\item<6-> Solving the system of equations with the dimension $z$.
		\item<7-> If necessary, undo step 1.
	\end{enumerate}
    
    \speech{KnotenMaschenanalyse}{1}{Um ein elektrisches Netzwerk vollständig zu berechnen, werden mehrere Arbeitsschritte empfohlen.}
    \speech{KnotenMaschenanalyse}{2}{Zuerst sollte das elektrische NEtzwerk so weit wie möglich vereinfacht werden. 
    Dass betrifft Beispielsweise in Reihe oder Parallel geschaltete Widerstände, welche zusammengefasst werden können.}
    \speech{KnotenMaschenanalyse}{3}{Folgend wird empfohlen, für eine bessere übersicht die Strompfeile und Spannungspfeile 
    in das elektrische Netzwerk zu zeichnen und einen zugehörigen Grafen zu erstellen.}
    \speech{KnotenMaschenanalyse}{4}{Hieraus können die K minus Eins Knotengleichungen}
    \speech{KnotenMaschenanalyse}{5}{sowie die Z minus K plus Eins linear unabhängigen Maschengleichungen aufgestellt werden.}
    \speech{KnotenMaschenanalyse}{6}{Das so entstandene Gleichungssystem mit der Dimension Z muss nun gelöst werden.}
    \speech{KnotenMaschenanalyse}{7}{Gegebenenfalls sollte je nach Bedarf die Vereinfachung des elektrischen Netzwerkes wieder Rückgängig 
    gemacht werden um eventuelle Teilspannungen oder Teilströme berechnen zu können.}
\end{frame}

\begin{frame} 
    \renewcommand{\figurename}{Figure}
	\ftx{Knoten- und Maschenanalyse}

	\s{
		The following example network is considered:
	}

	\fu{
		\input{Tikz/src/beispielanalyse.tex}
	}{\bf Example network for node and mesh analysis. \label{BildBeispielAnalyse}}

    \speech{KnotMaschNetzwerk}{1}{Folgend wird gezeigt, wie die Gleichungen zu einem elektrischen Netzwerk erstellt werden. 
    In diesem Netzwerk sind die Strompfeile bereits verzeichnet.}
\end{frame}

\begin{frame}
    \renewcommand{\figurename}{Figure} 
	\ftx{Knoten- und Maschenanalyse}

	\s{
		The corresponding graph and the meshes defined by it are shown in Figure \ref{BildKnotenMaschen}.
        A different arrangement of the meshes is possible. As shown, meshes can also extend beyond branches. 
	}

	\fu{
		\input{Tikz/src/knotenmaschen.tex}
		}{{\bf Graph with mesh, for example.} The tree branches are coloured red and the mesh blue. \label{BildKnotenMaschen}}
    
    \speech{KnotMaschGraph}{1}{Für das Netzwerk wurde im zugehörigen Grafen die vorgestellte Baumstruktur gewählt. 
    Bei den Maschen wird außer bei der Masche M 2 kein Zweig gekreuzt. 
    Die Wahl der Maschen und der Baumstruktur kann auch anders erfolgen.}
\end{frame}

\begin{frame} 
	\ftx{Knoten- und Maschenanalyse}
	
	\s{
		The network has $z=8$ branches and $k=5$ nodes. Therefore, $k-1=4$ node equations and $z-k+1=4$ 
        mesh equations must be set up. Node 5 has four branches, the others only three. Therefore, the 
        node equation for node 5 is omitted.

        \begin{align*}
			K_1:\qquad& I_1+I_4-I_5=0\\
			K_2:\qquad& I_2-I_4-I_8=0\\
			K_3:\qquad& -I_1-I_3+I_6=0\\
			K_4:\qquad& I_3-I_7+I_8=0\\
			M_1:\qquad& R_2I_2 + R_4I_4 + R_5I_5 - U_{2} = 0\\
			M_2:\qquad& R_4I_4 + R_5I_5 - R_7I_7 - R_8I_8 = 0\\
			M_3:\qquad& -R_1I_1 - R_5I_5 - R_6I_6 + U_{1} = 0\\
			M_4:\qquad& R_3I_3 + R_6I_6 + R_7I_7 - U_{3} = 0
		\end{align*}
	}
	\b{
		\begin{align*}
			\only<1>{K_1:\qquad& I_1+I_4-I_5=0\\
			K_2:\qquad& I_2-I_4-I_8=0\\
			K_3:\qquad& -I_1-I_3+I_6=0\\
			K_4:\qquad& I_3-I_7+I_8=0\\}
			\only<2>{M_1:\qquad& R_2I_2 + R_4I_4 + R_5I_5 - U_{2} = 0\\
			M_2:\qquad& R_4I_4 + R_5I_5 - R_7I_7 - R_8I_8 = 0\\
			M_3:\qquad& -R_1I_1 - R_5I_5 - R_6I_6 + U_{1} = 0\\
			M_4:\qquad& R_3I_3 + R_6I_6 + R_7I_7 - U_{3} = 0}
		\end{align*}
    }

    \speech{KnotMaschGleichung}{1}{Es müssen K minus 1 Knotengleichungen aufgestellt werden. 
    In dem Netzwerk sind K gleich 5 Knoten. Somit werden vier Knotengleichungen aufgestellt.
    Hier werden die Gleichungen für die Knoten K 1, K 2, K 3 und K 4 aufgestellt. 
    In den Knoten K 1 fließen die Ströme I 1 und I 4 hinein und der Strom I 5 hinaus.
    %In den Knoten K 2 fließen der Strom I 2 hinein und die Ströme I 4 und I 8 hinaus.
    %In den Knoten K 3 fließen die Ströme I 1 und I 3 hinaus und der Strom I 6 hinein.
    %In den Knoten K 4 fließen die Ströme I 3 und I 8 hinein und der Strom I 7 hinaus.
    Zur Erinnerung: Die in einen Knoten fließenden Ströme werden mit einem positiven Vorzeichen
    und die herausfließenden Ströme mit einem negativen Vorzeichen versehen.}
    \speech{KnotMaschGleichung}{2}{Die eingezeichneten Maschen verlaufen allesamt rechtsdrehend im Uhrzeigersinn.
    Die Spannungspfeile über den elektrischen Widerständen ergeben sich aus den zugehörigen Strompfeilen.
    Damit sind für die Masche M 1 die Spannungspfeile der Widerstände alle übereinstimmend mit dem Maschenumlauf.
    Nur die Spannung der Spannungsquelle U 2 ist dem Maschenumlauf entgegengesetzt und wird deswegen negativ notiert. 
    Zur Erinnerung: Verlaufen der Spannungspfeil und die Maschenrichtung in die selbe Richtung wird ein positives Vorzeichen benutzt.
    Für den Fall, dass die Richtungen entgegengesetzt sind, wird ein negatives Vorzeichen verwendet.}
\end{frame}

\begin{frame} 
	\ftx{Knoten- und Maschenanalyse}

	\s{
		To make the system of equations easier to solve and to increase clarity, it can be represented in 
        matrix notation. With a little practice, matrix notation can also be used immediately.
	}

		\begin{equation*}	
			\left[\begin{array}{cccccccc}
			1&0&0&1&-1&0&0&0\\
			0&1&0&-1&0&0&0&-1\\
			-1&0&-1&0&0&1&0&0\\
			0&0&1&0&0&0&-1&1\\
			0&R_2&0&R_4&R_5&0&0&0\\
			0&0&0&R_4&R_5&0&-R_7&-R_8\\
			-R_1&0&0&0&-R_5&-R_6&0&0\\
			0&0&R_3&0&0&R_6&R_7&0
			\end{array}\right] 
			\cdot 
			\left(\begin{array}{c}
			I_1\\I_2\\I_3\\I_4\\I_5\\I_6\\I_7\\I_8
			\end{array}\right) 
			= 
			\left(\begin{array}{c} 
			0\\0\\0\\0\\U_{2}\\0\\-U_{1}\\U_{3}
			\end{array}\right)
		\end{equation*}
		
	\s{
		This 8th order system of equations can now in principle be solved using the methods known from mathematics. 
        However, we do not want to do that here, as simpler methods are presented below that are easier to calculate.
	}
		
    \speech{KnotMaschMatrix}{1}{Die Knotengleichungen und Maschengleichungen lassen sich in einer Matrix zusammenfassen.}
\end{frame}




\newpage
%Beispiel
\begin{frame}
    \renewcommand{\figurename}{Figure}
    \ftx{Beispiel Netzwerkanalyse}
	\begin{expl}{Network analysis}{}

    \s{
        The electrical network shown in Figure \ref{BildBeispielaufgabe} is given. The following 
        tasks are to be completed:
    }

    \begin{enumerate}
       
        \only<1>{		
            \item[] 
            \b{
                Gegeben ist das angegebene elektrische Netzwerk. Die folgenden Aufgaben sollen bearbeitet werden: 
            }
            
            \begin{enumerate}
                \item[a)] Plotting currents and voltages
                \item[b)] Setting up the node equations
                \item[c)] Setting up the mesh equation $M_1$
                \item[d)] Setting up the mesh equation $M_2$ without the expression $U_1$
            \end{enumerate}

            \fu{
                \input{Tikz/src/beispielaufgabe1.tex}
            }{{\bf Example.} Network analysis using Kirchhoff's laws. \label{BildBeispielaufgabe}}
            }

			\only<2>{			
				\item[a)] Specification of currents and voltages:
				}
			\only<2-5>{			
				\item[] 
                \begin{figure}[H]
                    \resizebox{0.65\textwidth}{!}{%  
                        \input{Tikz/src/beispielaufgabe2.tex}
                        }\centering                                                           
                    \end{figure}
				}
			\only<3>{
				\item[b)] Knot equations:
				\onslide<3>{
					\begin{eqa}
						K_1: I_0-I_1-I_2-I_3&=0 \nonumber\\
                        K_2: -I_0+I_1+I_2+I_3&=0 \nonumber
					\end{eqa}
				}
			}
			\only<4>{
				\item[c)] Setting up the mesh equation $M_1$ and $M_2$:
				\onslide<4>{
					\begin{eqa}
						M_1: U_1+U_\mathrm{R2}-U_\mathrm{R1}&=0 \nonumber\\
                        M_2: U_2+U_\mathrm{R3}-U_\mathrm{R2}-U_1&=0 \nonumber
				    \end{eqa}
				}
			}
            \only<5>{
				\item[d)] Setting up the mesh equation $M_2$ without the expression $U_1$
				\onslide<5>{
					\begin{eqa}
						M_1: U_1&=U_\mathrm{R1}-U_\mathrm{R2} \nonumber\\
                        M_2: U_2+U_\mathrm{R3}-U_\mathrm{R2}-(U_\mathrm{R1}-U_\mathrm{R2})&=0 \nonumber\\
                        M_2: U_2+U_\mathrm{R3}-\cancel{U_\mathrm{R2}}-U_\mathrm{R1}+\cancel{U_\mathrm{R2}}&=0 \nonumber\\
                        M_2: U_2+U_\mathrm{R3}-U_\mathrm{R1}&=0 \nonumber
				    \end{eqa}
				}
			}
		\end{enumerate}
	
    \speech{KnotMaschBsp}{1}{Als Beispiel wird nun das dargestellte elektrische Netzwerk analysiert.
    Das Netzwerk enthält die folgenden Komponenten: Die Stromquelle I 0. Die Sannungsquellen U 1 und U 2.
    Sowie die Widertstände R 1, R 2 und R 3.
    Hierfür sollen die Ströme und Spannungen der Widerstände eingezeichnet werden.
    Danach sollen die Knotengleichungen und die Maschengleichungen aufgestellt werden.
    Hierbei soll für die Masche M 2 der Ausdruck der Spannungsquelle U 1 vermieden werden.}
    \speech{KnotMaschBsp}{2}{Hier werden sämtliche Strompfeile und Spannungspfeile dargestellt. 
    Außerdem wurden die Knoten K 1 und K 2 definiert.}
    \speech{KnotMaschBsp}{3}{Als nächstes werden die Knotengleichungen aufgestellt.
    In den Knoten K 1 fließt der Strom I Null. Die übrigen Ströme fließen aus dem Knoten heraus.
    Beim Knoten K 2 fließt der Strom I Null aus dem Knoten heraus und die übrigen Ströme fließen in den Knoten hinein.}
    \speech{KnotMaschBsp}{4}{Folgend werden die Maschengleichungen aufgestellt.
    Bei der Masche M 1 verlaufen die Spannungen U 1 und U R 2 in der Richtung des Maschenverlaufes. 
    Die Spannung U R 1 ist dem Maschenumlauf entgegengesetzt.
    Bei der Masche M 2 sind die Spannungen U 2 und U R 3 im Verlauf der Masche, die Spannungen U R 2 und U 1 verlaufen entgegen des Verlaufes der Masche.}
    \speech{KnotMaschBsp}{5}{Nun soll bei der Masche der Ausdruck U 1 vermieden werden. 
    Durch eine Umformulierung der Maschengleichung M 1 bekommen wir einen Ausdruck für U 1, welchen wir in die Maschengleichung von M 2 einsetzen können.
    Durch das Ausklammern und Aufsummieren erhalten wir so eine MAschengleichung für die MAsche M 2 in welcher der Ausdruck von U 1 vermieden wird.}
    \end{expl}

    
\end{frame}







